\section{Preliminaries}\label{sec:preliminaries}

% ~0.4 page total.
% Notation (agreed 2026-09-02): P transition kernel, H episode horizon,
% \mathcal{T} task space, \rho_0 initial state distribution, g^\star target goal,
% \epsilon goal tolerance.
% When porting the pseudocode: p_0 -> \rho_0, G -> g^\star, T_B unchanged.

\subsection{Goal-Conditioned Markov Decision Process}\label{sec:problem_formulation}

We model the task as a goal-augmented Markov decision process (GA-MDP)~\citep{liu2022goal} with state space $\mathcal{S}$, action space $\mathcal{A}$, transition kernel $P$, discount factor $\gamma$, initial state distribution $\rho_0$, and goal space $\mathcal{G}$.
A known mapping $\phi: \mathcal{S} \to \mathcal{G}$ projects a state onto the goal it achieves~\citep{andrychowicz2017hindsight}, and a goal $g$ is reached in state $s$ if $\lVert \phi(s) - g \rVert \le \epsilon$ for a tolerance $\epsilon$.
The reward is sparse and binary: the agent receives $+1$ in every step in which the goal is reached and $0$ otherwise. Episodes have a fixed horizon $H$ and do not terminate when the goal is reached.

The agent acts with a goal-conditioned policy $\pi_\theta(a \mid s, g)$~\citep{schaul2015universal}, trained with PPO~\citep{schulman2017proximal}.
%  in massively parallel simulation~\citep{rudin2022learning}.
We denote by $R(\tau, g)$ the discounted return of a rollout $\tau$ under goal $g$. Because the reward keeps accruing at the goal, the return measures how quickly and how reliably the policy reaches the goal and stays there, rather than a binary success.
The \emph{target task} is to reach a single target goal $g^\star \in \mathcal{G}$ from the initial state distribution $\rho_0$, and the learning objective is to maximize the expected return on this task.

\subsection{Curriculum}\label{sec:curriculum}

Optimizing the target task directly fails under sparse rewards and long horizons: a randomly initialized policy essentially never reaches $g^\star$ from $\rho_0$ and thus receives no learning signal~\citep{florensa2017reverse}.
Curriculum learning instead trains on a sequence of easier tasks~\citep{narvekar2020curriculum, portelas2020automatic}. In our setting, a task is a start--goal pair, and a \emph{curriculum} decides in every training iteration which pairs the policy is trained on, based on the experience collected with the current policy~\citep{tao2024reverse}. The intermediate tasks are auxiliary, and the objective remains the target task.
A \emph{start-state curriculum} keeps the goal fixed at $g^\star$ and proposes intermediate start states~\citep{florensa2017reverse, tao2024reverse}, whereas a \emph{goal curriculum} keeps the start states drawn from $\rho_0$ and proposes intermediate goals~\citep{florensa2018automatic, ren2019exploration}. Our method does both at once: it trains one shared goal-conditioned policy on intermediate starts and intermediate goals, with the training budget split evenly between the two (Section~\ref{sec:method}).

\subsection{Assumptions}\label{sec:assumptions}

Analogous to \citet{florensa2017reverse}, we assume that the simulator can be reset to arbitrary states and simulated from them under random actions, that the initial distribution $\rho_0$ can be sampled, that one goal state $s^{g} \in \mathcal{S}$ with $\phi(s^{g}) = g^\star$ is known, and that a bounded task space $\mathcal{T}$ that contains the target task is known and can be sampled uniformly.
